%
% $Id: slides.tex 4228 2010-11-24 17:55:12Z pcoca $
%
%
% Compilar a .pdf con LaTeX (pdflatex)
% Es necesario instalar Beamer (paquete latex-beamer en Debian)
%

%
% Gráficos:
% Los gráficos pueden suministrarse en PNG, JPG, TIF, PDF, MPS
% Los EPS deben convertirse a PDF (usar epstopdf)
%

\documentclass{beamer}
\usetheme{Warsaw}
\usebackgroundtemplate{\includegraphics[width=\paperwidth]{format/libresoft-bg.png}}
\usepackage[spanish]{babel}
\usepackage[utf8]{inputenc}
\usepackage{graphics}
\usepackage{amssymb} % Simbolos matematicos
\usepackage{url}

%\definecolor{libresoftgreen}{RGB}{162,190,43}
%\definecolor{libresoftblue}{RGB}{0,98,143}

%\setbeamercolor{titlelike}{bg=libresoftgreen}

%% Metadatos del PDF.
\hypersetup{
  pdftitle={QSOS: Qualification and Selection of Software Open Source},
  pdfauthor={F. Ortega, D. Izquierdo, P. Coca},
  pdfcreator={GSyC/Libresoft},
  pdfproducer=PDFLaTeX,
  pdfsubject={nn},
}
%%


\AtBeginSection[]
{
  \begin{frame}<presentation>
    \frametitle{Index}
    \tableofcontents[current]
  \end{frame}
}


\begin{document}

\title{QSOS}
\subtitle{Qualification and Selection of Software Open Source}
\institute{\\pcoca@libresoft.es\\
GSyC/Libresoft}
\author{Felipe Ortega, Daniel Izquierdo, Pedro Coca}
\date{\today}

\frame{
\maketitle
\begin{center}
\includegraphics[width=6cm]{format/gsyc-urjc}
\end{center}
}


% Si el titulo o el autor se quieren acortar para los pies de página
% se pueden redefinir aquí:
%\title{Titulo corto}
%\author{Autores abreviado}


%% LICENCIA DE REDISTRIBUCION DE LAS TRANSPAS
\frame{
~
\vspace{4cm}

\begin{flushright}
{\tiny
(cc) 2010 Felipe Ortega, Daniel Izquierdo, Pedro Coca. \\
Some rights reserved. This document is distributed under the Creative \\
            Commons Attribution-ShareAlike 3.0 licence, available in \\
            http://creativecommons.org/licenses/by-sa/3.0/

%  Este documento (o uno muy similar) está disponible en \\
%  \url{http://gsyc.escet.urjc.es/~jjamor/}
}
\end{flushright}
}
%%

%%%%%%
%Transpas separadas por \begin{frame}
%%%%%%%%%%%%%%%%%%%%%%%%\end{frame}

\section{Introduction}

\begin{frame}
\frametitle{What is QSOS?}
\begin{itemize}
\item QSOS Manifesto
\begin{itemize}
 \item Analyse needs and limitations in software adoption.
 \item Evaluate functional and technical requirements.
 \item Formalise a methodology.
 \item Results reuse.
 \item Provide a free method.
\end{itemize}

\item Raphaël Semeteys at Atos Origin release in 2004 the first official QSOS version: 1.3.
\item Afert Re-licensing, aditions and tranlations last released was done in 2006. QSOS 1.6. 
\item GNU Free Documentation License
\end{itemize}
\end{frame}

%%%%%%%%%%%%%%%%%%%%%%%%%%%%%%%%%%%%%%%%%%%%%%%%%%%%%%%%%%%%%%

\begin{frame}
 \frametitle{Target Audience of QSOS}
 \begin{itemize}
 \item Standard Software Quality Methodologies audience.
 \item + FLOSS Communities / Actors.
 \end{itemize}
\end{frame}


%%%%%%%%%%%%%%%%%%%%%%%%%%%%%%%%%%%%%%%%%%%%%%%%%%%%%%%%%%%%%%

\section{QSOS General Process}

%%%%%%%%%%%%%%%%%%%%%%%%%%%%%%%%%%%%%%%%%%%%%%%%%%%%%%%%%%%%%%

\begin{frame}
\frametitle{QSOS Process overview}
 \begin{itemize}
 \item Four independents steps
   \begin{itemize}
   \item Definition.
   \item Evaluation.
   \item Qualification.
   \item Selection.
   \end{itemize}
 \item Iterative process
 \item Tools provided
   \begin{itemize}
   \item O3S: Open Source Selection Software.
   \end{itemize}

 \end{itemize}
\end{frame}

%%%%%%%%%%%%%%%%%%%%%%%%%%%%%%%%%%%%%%%%%%%%%%%%%%%%%%%%%%%%%%

\begin{frame}
\frametitle{QSOS Process overview}
\begin{center}
\begin{figure}
 \includegraphics[height=6cm]{figs/QSOS_Process_Overview.png}
\end{figure}
\end{center}
\end{frame}

%%%%%%%%%%%%%%%%%%%%%%%%%%%%%%%%%%%%%%%%%%%%%%%%%%%%%%%%%%%%%%

\begin{frame}
 \frametitle{QSOS Process: Definition}
 \begin{itemize}
 \item Initial stage of the process.
 \item Define Frame of references.
   \begin{itemize}
   \item Software Families.
   \item Types of licenses.
      \begin{itemize}
      \item Ownership.
      \item Virality.
      \item Inheritance.
      \end{itemize}
   \item Types of communities.
      \begin{itemize}
      \item Insulated developer.
      \item Group of developers.
      \item Organization of developers.
      \item Legal entity.
      \item Commercial entity.
      \end{itemize}
   \end{itemize}
 \end{itemize}
\end{frame}

%%%%%%%%%%%%%%%%%%%%%%%%%%%%%%%%%%%%%%%%%%%%%%%%%%%%%%%%%%%%%%

\begin{frame}
\frametitle{QSOS Process: Definition}
\begin{center}
\begin{figure}
 \includegraphics[height=6cm]{figs/QSOS_Process_Definition.png}
\end{figure}
\end{center}
\end{frame}

%%%%%%%%%%%%%%%%%%%%%%%%%%%%%%%%%%%%%%%%%%%%%%%%%%%%%%%%%%%%%%

\begin{frame}
 \frametitle{QSOS Process: Evaluation}
 \begin{itemize}
 \item Based on the Community information.
 \item Collected the relevant information.
 \item Targets:
    \begin{itemize}
     \item Identity Card (Not directly scored!)
     \item Evaluation Sheet
    \end{itemize}
 \end{itemize}
\end{frame}

%%%%%%%%%%%%%%%%%%%%%%%%%%%%%%%%%%%%%%%%%%%%%%%%%%%%%%%%%%%%%%

\begin{frame}
 \frametitle{QSOS Process: Evaluation}
 \begin{itemize}
 \item Identity Card.
    \begin{itemize}
     \item General Information
     \item Existing Services
     \item Functional and technical aspects
     \item Synthesis
    \end{itemize}

 \item Evaluation Sheet.
    \begin{itemize}
     \item Functional Coverage
     \item Risks (User's perspective)
     \item Risks (Service provider's perspective)
    \end{itemize}

 \end{itemize}
\end{frame}

%%%%%%%%%%%%%%%%%%%%%%%%%%%%%%%%%%%%%%%%%%%%%%%%%%%%%%%%%%%%%%

\begin{frame}
 \frametitle{QSOS Process: Evaluation Sheet}
 \begin{itemize}
 \item Scoring
    \begin{itemize}
     \item Criteria: [0-2]
    \end{itemize}

 \item Functional Coverage
    \begin{itemize}
     \item Not Covered(0)/Partially Covered(1)/Completely Covered(2)
    \end{itemize}
 \item Risks (User's perspective)
    \begin{itemize}
     \item Intrinsic durability (Maturity, Adoption, etc)
     \item Industrialised solution (Services, Documentation, etc)
     \item Integration (Packaging, etc)
     \item Technical adaptability (Modularity, code modification, etc)
     \item Strategy (License, Copyright, Independence, etc)
    \end{itemize}
 \item Risks (Service provider's perspective)
    \begin{itemize}
     \item Services providing (Maintenability, Code expertise, etc)
    \end{itemize}

 \end{itemize}
\end{frame}

%%%%%%%%%%%%%%%%%%%%%%%%%%%%%%%%%%%%%%%%%%%%%%%%%%%%%%%%%%%%%%

\begin{frame}
\frametitle{QSOS Process: Evaluation}
\begin{center}
\begin{figure}
 \includegraphics[height=6cm]{figs/QSOS_Process_Evaluation.png}
\end{figure}
\end{center}
\end{frame}

%%%%%%%%%%%%%%%%%%%%%%%%%%%%%%%%%%%%%%%%%%%%%%%%%%%%%%%%%%%%%%

\begin{frame}
 \frametitle{QSOS Process: Qualification}
 \begin{itemize}
 \item Define filters.
    \begin{itemize}
     \item Take into account needs and limitations
     \item Specific context!
    \end{itemize}

 \item Filter on ID card.
 \item Filter on Functional grid.
    \begin{itemize}
     \item Required functionality
     \item Optional functionality
     \item Not required functionality
    \end{itemize}

 \end{itemize}
\end{frame}

%%%%%%%%%%%%%%%%%%%%%%%%%%%%%%%%%%%%%%%%%%%%%%%%%%%%%%%%%%%%%%


\begin{frame}
\frametitle{QSOS Process: Qualification}
\begin{center}
\begin{figure}
 \includegraphics[height=5cm]{figs/QSOS_Process_Qualification.png}
\end{figure}
\end{center}
\end{frame}

%%%%%%%%%%%%%%%%%%%%%%%%%%%%%%%%%%%%%%%%%%%%%%%%%%%%%%%%%%%%%%

\begin{frame}
 \frametitle{QSOS Process: Selection}
 \begin{itemize}
 \item Identify software meeting the needs/requirements
 \item Compare sofware, targeting selection.
 \item Selection modes.
    \begin{itemize}
     \item Strict selection
     \item loose selection
    \end{itemize}

 \end{itemize}
\end{frame}
%%%%%%%%%%%%%%%%%%%%%%%%%%%%%%%%%%%%%%%%%%%%%%%%%%%%%%%%%%%%%%


\begin{frame}
\frametitle{QSOS Process overview}
\begin{center}
\begin{figure}
 \includegraphics[height=5cm]{figs/QSOS_Process_Selection.png}
\end{figure}
\end{center}
\end{frame}

%%%%%%%%%%%%%%%%%%%%%%%%%%%%%%%%%%%%%%%%%%%%%%%%%%%%%%%%%%%%%%

\section{References}

\begin{frame}
 \frametitle{References}
 \begin{itemize}
  \item QSOS Website. \textit{http://www.qsos.org}.
  \item O3S Tool \textit{http://www.qsos.org/o3s/}
  \item QSOS Community Web Site \textit{https://savannah.nongnu.org/projects/qsos/}
  \end{itemize}

\end{frame}

%%%%%%%%%%%%%%%%%%%%%%%%%%%%%%%%%%%%%%%%%%%%%%%%%%%%%%%%%%%%%%

\end{document}
